% ============================================================
\section{数据类型与类型系统}
\subsection*{Data Types and Typing}

本节说明 PXL 支持的数据类型，以及决定 PXL 内类型如何工作的因素。

变量可显式声明类型。例如：
\begin{lstlisting}[style=pxl]
i : integer = 1 ;           // i is declared as an integer
\end{lstlisting}

变量也可根据运算符的返回类型隐式定型。例如：
\begin{lstlisting}[style=pxl]
a : double = 1.1203;
b : double = 2.0213;
c = a + b;                     // c is implicitly the type of "double"
\end{lstlisting}

\subsection{受支持的数据类型}
\subsubsection*{Supported Data Types}

PXL 支持多数基本类型以及一些工具相关数据类型。

\subsubsection{基本类型}
\paragraph*{Primitive Types}

表~\ref{tab:pxl-primitive} 给出 PXL 支持的基本数据类型。

\begin{longtable}{@{}>{\raggedright\arraybackslash\ttfamily}p{0.34\textwidth} p{0.58\textwidth}@{}}
\caption{基本类型（Primitive Types）}\label{tab:pxl-primitive}\\
\toprule
\textbf{\textrm{数据类型}} & \textbf{\textrm{说明}} \\
\midrule
\endfirsthead
\toprule
\textbf{\textrm{数据类型}} & \textbf{\textrm{说明}} \\
\midrule
\endhead
\bottomrule
\endfoot
integer & 整数；大小由实现定义，一般为 32 位有符号值。 \\
double & 双精度浮点；大小在实现时定义，一般为 64 位值。 \\
boolean & true 或 false \\
string & 字符序列；支持用转义序列包含特殊字符。 \\
handle & 不透明数据类型，用于跟踪 PXL 程序本身之外的数据对象。 \\
constraint of double 与 constraint of integer & 连续 double 值或 integer 值集合的简洁表示。 \\
predicate\_handle & 布尔表达式，并对将层标识符求值为空/非空有特殊处理。 \\
\end{longtable}

\paragraph{integer}

整数声明示例如下：
\begin{lstlisting}[style=pxl]
i : integer;     // declared but not initialized
i : integer = 1; // declared and initialized
\end{lstlisting}

\paragraph{string}

定义零个或多个字符的字符串。\cmd{string} 类型由字符串字面量初始化。例如：
\begin{lstlisting}[style=pxl]
s : string = "xyzzy";
\end{lstlisting}

\cmd{string} 类型对任意 \cmd{string} 类型表达式支持表~\ref{tab:pxl-string-ops} 所示操作。

\begin{longtable}{@{}>{\raggedright\arraybackslash\ttfamily}p{0.18\textwidth} p{0.12\textwidth} p{0.12\textwidth} p{0.14\textwidth} p{0.34\textwidth}@{}}
\caption{字符串类型操作（String Type Operations）}\label{tab:pxl-string-ops}\\
\toprule
\textbf{\textrm{操作}} & \textbf{\textrm{e2 类型}} & \textbf{\textrm{返回}} & \textbf{\textrm{修改 e1}} & \textbf{\textrm{说明}} \\
\midrule
\endfirsthead
\toprule
\textbf{\textrm{操作}} & \textbf{\textrm{e2 类型}} & \textbf{\textrm{返回}} & \textbf{\textrm{修改 e1}} & \textbf{\textrm{说明}} \\
\midrule
\endhead
\bottomrule
\endfoot
e1.empty() & -- & boolean & 否 & 若字符串为空（即 \cmd{""}）则返回 true；否则返回 false。 \\
e1.size() & -- & integer & 否 & 返回字符串中的字符数。 \\
\end{longtable}

\paragraph{handle}

外部管理对象的内部表示。类型 \cmd{handle} 具有对每个外部管理对象唯一的字符串表示。

\paragraph{constraint of double 与 constraint of integer}

\cmd{constraint of double} 与 \cmd{constraint of integer} 是尺寸规则值的数学区间。这些约束为描述 DRC 间距距离等参数提供直观机制。

若未指定约束类型，则隐含该区间为 \cmd{constraint of double}。

约束数据类型对属于该类型的表达式支持表~\ref{tab:pxl-constraint-ops} 所示操作。

\begin{longtable}{@{}>{\raggedright\arraybackslash\ttfamily}p{0.20\textwidth} p{0.12\textwidth} p{0.14\textwidth} p{0.14\textwidth} p{0.30\textwidth}@{}}
\caption{constraint of double 与 constraint of integer 的操作}\label{tab:pxl-constraint-ops}\\
\toprule
\textbf{\textrm{操作}} & \textbf{\textrm{e2 类型}} & \textbf{\textrm{返回}} & \textbf{\textrm{修改 e1}} & \textbf{\textrm{说明}} \\
\midrule
\endfirsthead
\toprule
\textbf{\textrm{操作}} & \textbf{\textrm{e2 类型}} & \textbf{\textrm{返回}} & \textbf{\textrm{修改 e1}} & \textbf{\textrm{说明}} \\
\midrule
\endhead
\bottomrule
\endfoot
e1.contains(e2) & item & boolean & 否 & 若 e2 值落在 e1 约束指定的区间内则返回 true。 \\
e1.category() & -- & enum & 否 & 返回该约束在预定义枚举 constraint\_category 中的枚举元。见表~\ref{tab:pxl-constraint-cat}。 \\
e1.lo() & -- & integer/double & 否 & 对二元约束，返回下界：constraint of double 为 double；constraint of integer 为 integer。\textsuperscript{1} \\
e1.hi() & -- & integer/double & 否 & 对二元约束，返回上界：constraint of double 为 double；constraint of integer 为 integer。\textsuperscript{1} \\
\end{longtable}

\noindent\small
\textsuperscript{1}对一元约束（如 \cmd{==} 与 \cmd{>}），\cmd{e1.lo()} 与 \cmd{e1.hi()} 均返回该单一界限。

约束数据类型只能通过从另一同类型约束赋值，或从字面常量创建，如表~\ref{tab:pxl-constraint-cat} 所示。

\begin{longtable}{@{}>{\raggedright\arraybackslash\ttfamily}p{0.14\textwidth} p{0.48\textwidth} >{\raggedright\arraybackslash\ttfamily}p{0.28\textwidth}@{}}
\caption{约束类别（Constraint Categories）}\label{tab:pxl-constraint-cat}\\
\toprule
\textbf{\textrm{字面量}} & \textbf{\textrm{说明}} & \textbf{\textrm{约束类别}} \\
\midrule
\endfirsthead
\toprule
\textbf{\textrm{字面量}} & \textbf{\textrm{说明}} & \textbf{\textrm{约束类别}} \\
\midrule
\endhead
\bottomrule
\endfoot
[a,b] & 从 a 到 b 的闭区间，包含两端点 & CONSTRAINT\_GELE \\
(a,b) & 从 a 到 b 的开区间，不含两端点 & CONSTRAINT\_GTLT \\
[a,b) & 从 a 到 b 的半开区间，含 a 不含 b & CONSTRAINT\_GELT \\
(a,b] & 从 a 到 b 的半开区间，不含 a 含 b & CONSTRAINT\_GTLE \\
>= a & 大于等于 a 的全部数 & CONSTRAINT\_GE \\
> a & 大于 a 的全部数 & CONSTRAINT\_GT \\
<= a & 小于等于 a 的全部数 & CONSTRAINT\_LE \\
< a & 小于 a 的全部数 & CONSTRAINT\_LT \\
== a & 仅包含数 a 的区间 & CONSTRAINT\_EQ \\
!= a & 包含除 a 以外全部数的区间 & CONSTRAINT\_NE \\
\end{longtable}

\cmd{lo()} 与 \cmd{hi()} 对所有约束均有效。对单端约束，两函数均定义为返回该单一端点 \cmd{a}。对双端约束，\cmd{lo()} 返回 \cmd{a}，\cmd{hi()} 返回 \cmd{b}。

\paragraph{示例}
\textit{Examples}

示例如下。
\begin{itemize}
  \item Constraint of integer
  \begin{itemize}
    \item \cmd{[2,8]} 包含值 2、3、4、5、6、7、8
    \item \cmd{(2,8)} 包含值 3、4、5、6、7
  \end{itemize}
  \item Constraint of double
  \begin{itemize}
    \item \cmd{[2,8]} 包含该区间内全部实数（含 2.0 与 8.0）
    \item \cmd{(2,8)} 包含除 2.0 与 8.0 外该区间内的全部实数
  \end{itemize}
  \item 声明示例
\begin{lstlisting}[style=pxl]
a_value : constraint of double = [2.0,8.0];
\end{lstlisting}
  \item 赋值示例
\begin{lstlisting}[style=pxl]
distance <= 2.0                 // All values less than or equal to 2.0
distance = <= 2.0               // All values less than or equal to 2.0
distance = [ 2.0, 8.0 )         // All values >= 2.0 and < 8.0
count = 5                       // Count is exactly equal to 5
count = ==5                     // Count is exactly equal to 5
\end{lstlisting}
\begin{noteBox}
诸如 \cmd{==} 与 \cmd{<=} 的符号是约束的一部分。对期望约束类型的变量，赋值运算符 \cmd{=} 可选。但若右侧是约束类型变量，则必须使用赋值运算符：
\end{noteBox}
\begin{lstlisting}[style=pxl]
my_constraint: constraint of double = < 5.0;
external1(layer1, distance = my_constraint );
external1(layer1, distance <5);                              // same as above
external1(layer1, distance = < 5);                           // same as above
\end{lstlisting}
\end{itemize}

\paragraph{predicate\_handle}

括在双方括号 \cmd{[[\ldots]]} 中的布尔表达式。允许任意包含至少一个层名或违例名的布尔逻辑。层名在谓词表达式内有特殊含义：运行时若该层有数据则求值为 \cmd{true}，若层为空则求值为 \cmd{false}。违例对象在存在违例时为 \cmd{true}，无违例时为 \cmd{false}。

在 runset 函数中，\cmd{predicate\_handle} 参数根据表达式求值控制给定函数的行为：
\begin{itemize}
  \item \cmd{true}：函数按正常方式运行。
  \item \cmd{false}：函数产生空输出层。
\end{itemize}

例如：
\begin{lstlisting}[style=pxl]
temp_layer = and(poly, diff, predicate = [[metal1]]);
\end{lstlisting}

若 \cmd{metal1} 为空，则跳过 \cmd{and()} 函数，且 \cmd{temp\_layer} 为空。

\subsection{用户定义结构体与类型}
\subsubsection*{User-Defined Structures and Types}

除基本类型外，PXL 还支持表~\ref{tab:pxl-user-types} 所示的结构体与类型。

\begin{longtable}{@{}>{\raggedright\arraybackslash\ttfamily}p{0.28\textwidth} p{0.64\textwidth}@{}}
\caption{用户定义类型（User-Defined Types）}\label{tab:pxl-user-types}\\
\toprule
\textbf{\textrm{用户定义类型}} & \textbf{\textrm{主要特征}} \\
\midrule
\endfirsthead
\toprule
\textbf{\textrm{用户定义类型}} & \textbf{\textrm{主要特征}} \\
\midrule
\endhead
\bottomrule
\endfoot
newtype... & 定义类型。 \\
enum of... & 定义枚举元素列表。 \\
list of... & 定义元素列表。 \\
hash of... & 将一种数据类型的元素列表转换为另一种数据类型。 \\
struct of... & 定义由多个变量组成的复合类型。 \\
function & 定义可传递给另一用户定义函数、并可像普通变量一样调用的用户定义函数。 \\
\end{longtable}

\subsubsection{newtype...}
\paragraph*{newtype...}

\cmd{newtype} 声明一种新变量类型。它并不定义变量。每种类型都是唯一的。

合法类型定义示例如下：
\begin{lstlisting}[style=pxl]
/*
 * define a new type, direction.
 */

direction: newtype enum of {
   up, down, left, right, all // trailing comma is allowed
};

/* define two distinct types */
T1: newtype integer;
T2: newtype integer;

/* define objects of the two types */
o1: T1;
o2: T2;
\end{lstlisting}

非法类型定义示例如下：
\begin{lstlisting}[style=pxl]
/* given the definitions above, o1 and o2 are of
 * different types, therefore they cannot be compared
 * and this is an error
 */

if (o1 == o2) { ... }
\end{lstlisting}

下面是定义并使用名为 \cmd{direction} 的新数据类型的示例。它有四个允许取值。
\begin{lstlisting}[style=pxl]
direction : newtype enum of {
   NORTH,
   SOUTH,
   EAST,
   WEST
};
compass : direction; // Declare compass as type "direction"
compass = EAST;       // Assign a legal value to "compass"
\end{lstlisting}

\subsubsection{enum of...}
\paragraph*{enum of...}

定义一组与其他类型不同的名称（枚举字面量）。PXL 编码标准要求枚举字面量为全大写字符。

例如：
\begin{lstlisting}[style=pxl]
t3: newtype enum of {
   X,
   Y,
   Z,               // Trailing comma is not needed; silently discarded.
};
\end{lstlisting}

另一示例：
\begin{lstlisting}[style=pxl]
direction : newtype enum of {
   NORTH, SOUTH, EAST, WEST
};

compass : direction;                 // Declare compass as type "direction"
compass = EAST;                      // Assign a legal value to "compass"
\end{lstlisting}

枚举可与其他枚举类型共享名称而无错误：
\begin{lstlisting}[style=pxl]
t1_e : newtype enum of {
   NONE, TOP, BOTTOM, ALL
};
\end{lstlisting}

% --- 第8章后半见下方继续（由另一任务追加或合并）---
